\section{Related Work}
\label{sec:related}

\subsection{Benchmark Saturation Studies}

The saturation of ML benchmarks has attracted growing attention. Bowman and Dahl~\cite{Bowman2021} argue that standard NLU benchmarks have largely saturated and call for principled retirement criteria --- but their analysis is qualitative; they identify the \textit{need} for criteria without providing a quantitative threshold.

Liao et al.~\cite{Liao2022} conduct the largest systematic study of benchmark saturation, analyzing near-saturation trends across 3,765 benchmarks using time-based metrics and CoV. Their analysis confirms saturation is a systemic phenomenon --- but their metric is time-based (years to saturation), not paper\_count-based, and targets aggregate trends rather than discrete structural breaks.

Akhtar et al.~\cite{Akhtar2026} introduce the $S_{\text{index}}$, a composite saturation metric for 60 LLM benchmarks, characterizing saturation as a continuous property. It does not apply change-point detection to PwC-internal CoV-vs-paper\_count data and does not produce a paper\_count threshold. Vasudevan et al.~\cite{Vasudevan2022} analyze ImageNet saturation in depth, documenting that top-1 accuracy above 90\% no longer discriminates between models --- providing independent evidence for the exploration-phase variance of pre-saturation benchmarks. Cao and Zhao~\cite{Cao2025} demonstrate that LLM benchmark scores are increasingly inflated by test-set pretraining, further motivating saturation detection.

\textbf{The gap:} No prior work applies change-point detection to PwC-internal CoV-vs-paper\_count data to produce a discrete, paper-count-based retirement threshold.

\subsection{Change-Point Detection Methods}

PELT (Pruned Exact Linear Time)~\cite{Killick2012} is an exact dynamic programming algorithm for detecting an unknown number of change-points in 1D signals with a penalized cost function. Wang, Lin and Willett~\cite{Wang2019} establish the VPWBS algorithm's $O_p(1/n)$ localization rate for regression change-points, providing theoretical grounding for PELT's performance at small sample sizes (our $N=115$). The \texttt{ruptures} Python library~\cite{Truong2020} provides a well-maintained PELT implementation (2,000+ GitHub stars).

Existing benchmark analysis methods treat saturation as a continuous property or apply domain-specific heuristics. We instead treat the CoV-vs-paper\_count series as a signal in which a structural break should be detected and statistically validated.

\subsection{Benchmark Evaluation and Leaderboard Dynamics}

Papers With Code~\cite{PwC2019} maintains a comprehensive leaderboard database making it a uniquely rich source for benchmark lifecycle analysis. Prior work using PwC data has examined research dynamics but not structural breaks in performance variability organized by publication count. CoV ($\sigma/\mu$ per benchmark) captures score diversity in a scale-invariant manner, suitable for pooling across heterogeneous metric scales~\cite{Liao2022}.

Our work is the first to apply PELT to PwC-internal residual CoV data, producing a quantitative $\text{paper\_count}^*$ threshold grounded in change-point theory.
